\setcounter{chapter}{2}
\setcounter{section}{0}
\counterwithout{section}{chapter}
\chapter{Valeurs et Opérateurs Booléens}

\section{Les Valeurs Booléennes : 0 et 1}

Les valeurs booléennes (du nom du mathématicien George Boole) sont la base de toute l'informatique numérique et de l'électronique. Elles ne peuvent prendre que deux états possibles :
\begin{itemize}
    \item 1 (Vrai / True) : Représente un état actif, une condition satisfaite, ou un circuit fermé (passage du courant).
    \item 0 (Faux / False) : Représente un état inactif, une condition non satisfaite, ou un circuit ouvert.
\end{itemize}
   

Une \motcle{variable booléenne} est une variable qui ne peut stocker que l'une de ces deux valeurs.

\section{Les Opérateurs Booléens Fondamentaux}

Il existe trois opérateurs booléens fondamentaux : \motcle{ET (AND), OU (OR) et NON (NOT).}

\subsection{ L'Opérateur NON (NOT)}

L'opérateur NON (ou négation) est l'opérateur le plus simple : il inverse la valeur booléenne sur laquelle il s'applique.

\begin{center}
\begin{tabular}{lllll}
\cline{1-2}
\multicolumn{1}{|l|}{A} & \multicolumn{1}{l|}{NOT A} &  &  &  \\ \cline{1-2}
\multicolumn{1}{|l|}{1} & \multicolumn{1}{c|}{0}     &  &  &  \\ \cline{1-2}
\multicolumn{1}{|l|}{0} & \multicolumn{1}{c|}{1}     &  &  &  \\ \cline{1-2}
                        &                            &  &  & 
\end{tabular}
\end{center}

\subsection{L'Opérateur ET (AND)}

L'opérateur ET (ou conjonction) est un opérateur à deux entrées (A et B). Le résultat est 1 seulement si les deux entrées sont 1. Dans tous les autres cas, le résultat est 0.

\begin{center}
\begin{tabular}{|l|l|c|ll}
\cline{1-3}
A & B & A AND B &  &  \\ \cline{1-3}
0 & 0 & 0       &  &  \\ \cline{1-3}
0 & 1 & 0       &  &  \\ \cline{1-3}
1 & 0 & 0       &  &  \\ \cline{1-3}
1 & 1 & 1       &  &  \\ \cline{1-3}
\end{tabular}
\end{center}

\subsection{L'Opérateur OU (OR)}

L'opérateur OU (ou disjonction) est également un opérateur à deux entrées. Le résultat est 1 si au moins une des entrées est 1. Le résultat n'est 0 que si les deux entrées sont 0. Il s'agit du OU inclusif.

\begin{center}
\begin{tabular}{|l|l|c|l}
\cline{1-3}
A & B & A OR B  \\ \cline{1-3}
0 & 0 & 0     \\ \cline{1-3}
0 & 1 & 1     \\ \cline{1-3}
1 & 0 & 1     \\ \cline{1-3}
1 & 1 & 1     \\ \cline{1-3}
\end{tabular}
\end{center}

\subsection{Expressions Booléennes et Tables de Vérité}

Une expression booléenne combine des variables booléennes avec des opérateurs booléens.

La table de vérité d'une expression booléenne liste toutes les combinaisons possibles des valeurs d'entrée et le résultat de l'expression pour chaque combinaison.\\

Exemple : Dressons la table de vérité de l'expression E = (NOT A) OR B.\\
\begin{itemize}
    \item Identifier les entrées : Il y a deux variables (A et B), donc il y 4 combinaisons possibles.
    \item  Identifier les étapes intermédiaires : L'opération la plus interne est NOT A.
    \item Calculer le résultat final : L'opération finale est le OR entre le résultat de l'étape 2 et B.
\end{itemize}
\begin{center}
\begin{tabular}{|l|l|l|l|l}
\cline{1-4}
A & B & NOT A & (NOT A) OR B &  \\ \cline{1-4}
0 & 0 & 1     & 1 OR 0 =   &  \\ \cline{1-4}
0 & 1 & 1     & 1 OR 1 =    &  \\ \cline{1-4}
1 & 0 & 0     & 0 OR 0 =    &  \\ \cline{1-4}
1 & 1 & 0     & 0 OR 1 =    &  \\ \cline{1-4}
\end{tabular}
\end{center}

\begin{Avous}
    Construisez la table de vérité de l'expression booléenne E = A AND (NOT B).
    \vspace{3.5cm}
\end{Avous}

\section{OU Exclusif}

Le OU Exclusif (souvent noté XOR ou $\oplus$) est un autre opérateur très important.

Le résultat est 1 si et seulement si les deux entrées sont différentes. Si elles sont identiques (0,0 ou 1,1), le résultat est 0.

\begin{Avous}
    Calculez la valeur de l'expression E=(1 AND 0) XOR (NOT 1)
    \begin{enumerate}
        \item 1 AND 0 =
        \item NOT 1 = 
        \item 
    \end{enumerate}
    \vspace{3cm}
\end{Avous}

\section{Application direct}
L'addition binaire sur un bit est un exemple parfait et direct d'application des opérateurs booléens, en particulier le XOR.

Quand on additionne deux bits (A et B), on obtient deux résultats :
\begin{itemize}
    \item La Somme : le bit de résultat.
   \item La Retenue : le bit à reporter à la colonne suivante.
\end{itemize}
\begin{center}
\begin{tabular}{|l|l|l|l|l}
\cline{1-4}
A & B & Somme & Retenue &  \\ \cline{1-4}
0 & 0 & 0     & 0       &  \\ \cline{1-4}
0 & 1 & 1     & 0       &  \\ \cline{1-4}
1 & 0 & 1     & 0       &  \\ \cline{1-4}
1 & 1 & 0     & 1       &  \\ \cline{1-4}
\end{tabular}
\end{center}
En observant le tableau , on voit que :
\begin{itemize}
    \item Somme (S) : Correspond exactement à l'opérateur 
    \item Retenue (R) : Correspond exactement à l'opérateur 
\end{itemize}

\begin{Avous}
    En utilisant les règles de l'addition binaire (A XOR B pour la somme, A AND B pour la retenue), effectuez l'addition binaire de 1012 et 0112.
    \vspace{5cm}
\end{Avous}

\section{Opérateurs Booléen dans un langage de programmation}
Dans de nombreux langages (Python, Java, C++, etc.), les opérateurs AND et OR sont évalués de manière séquentielle (ou à court-circuit : short-circuit evaluation).\\

\textbf{Dans une expression A AND B :}\\

Si l'expression A est Faux (0), alors le résultat final sera nécessairement Faux (0), quelle que soit la valeur de B. \textbf{Le programme n'évalue pas l'expression B}.

\begin{center}
    \begin{tabular}{|l|l|c|c|}
\hline
A        & B      & A and B & Evaluation de B ? \\ \hline
0 (faux) & 0 ou 1 & 0       & non               \\ \hline
1 (vrai) & 0 ou 1 & valeur de B & oui               \\ \hline
\end{tabular}
\end{center}

\begin{Avous}
Supposons qu'une variable X vaille 0. Dans l'expression E=(X!=0) AND (10 / X > 1). Est-ce qu'une division par zéro sera tentée dans un langage utilisant l'évaluation séquentielle ? Pourquoi ?\\
\vspace{3cm}
\end{Avous}

\newpage
\chapter*{Exercices Booléens}
% Absolument \! Voici trois petits exercices pour mettre en pratique les concepts que nous venons de voir.

% ## **3 Petits Exercices Booléens**
\begin{Exercise}[title = Évaluation d'Expressions]
Calculez la valeur (0 ou 1) de l'expression booléenne suivante en respectant l'ordre de priorité des opérateurs ($NOT$ \> $AND$ \> $OR$) :
 $$E = (1\ OR\ NOT\ (0\ AND\ 1))\ AND\ 0$$
\end{Exercise}

% **Réponse (Masquée) :**

% 1.  Priorité 1 (parenthèse interne $AND$ ) : $0\ AND\ 1 = 0$
% 2.  L'expression devient : $E = (1\ OR\ NOT\ 0)\ AND\ 0$
% 3.  Priorité 2 ($NOT$) : $NOT\ 0 = 1$
% 4.  L'expression devient : $E = (1\ OR\ 1)\ AND\ 0$
% 5.  Priorité 3 (parenthèse $OR$) : $1\ OR\ 1 = 1$
% 6.  L'expression devient : $E = 1\ AND\ 0$
% 7.  Priorité 4 ($AND$) : $1\ AND\ 0 = 0$

% **Résultat :** $E = 0$

% -----
\begin{Exercise}[title = Table de Vérité et Égalité]
\Question  Construisez la table de vérité pour les deux expressions $E_1$ et $E_2$. Les deux expressions sont-elles logiquement équivalentes ?

 $$E_1 = A\ XOR\ B$$
$$E_2 = (A\ OR\ B)\ AND\ (NOT\ (A\ AND\ B))$$

\begin{center}
    
\Large{
 \begin{tabular}{|l|l|l|l|l|l|l|}
    \hline
        A & B & A OR B & A AND B & NOT (A AND B) & E\_2 & E\_1 (A XOR B) \\ \hline
        & & & & & & \\ \hline
        & & & & & & \\ \hline
        & & & & & & \\ \hline
        & & & & & & \\ \hline
    \end{tabular}
}
\end{center}

\end{Exercise}
% ### **Exercice 2 : Table de Vérité et Égalité**

% 

% **Réponse (Masquée) :**

% | $A$ | $B$ | $A\ OR\ B$ | $A\ AND\ B$ | $NOT\ (A\ AND\ B)$ | $E_2$ | $E_1$ ($A\ XOR\ B$) |
% | :--: | :--: | :--------: | :---------: | :----------------: | :---: | :-----------------: |
% | $0$ | $0$ | $0$ | $0$ | $1$ | $0\ AND\ 1 = 0$ | $0$ |
% | $0$ | $1$ | $1$ | $0$ | $1$ | $1\ AND\ 1 = 1$ | $1$ |
% | $1$ | $0$ | $1$ | $0$ | $1$ | $1\ AND\ 1 = 1$ | $1$ |
% | $1$ | $1$ | $1$ | $1$ | $0$ | $1\ AND\ 0 = 0$ | $0$ |

% **Conclusion :** En comparant la colonne $E_1$ et la colonne $E_2$, on voit que les résultats sont identiques pour toutes les combinaisons.
% **Oui**, les deux expressions sont logiquement équivalentes. L'expression $E_2$ est d'ailleurs la définition du $XOR$ en utilisant uniquement les opérateurs fondamentaux $AND$, $OR$ et $NOT$.

% -----
\begin{Exercise}[title = Application Séquentielle]
Considérez le fragment de code (en pseudo-langage) suivant, où l'évaluation est \textbf{séquentielle} pour l'opérateur AND. Quel message sera affiché à l'écran si l'on exécute le code ?
\begin{lstlisting}
A = 1
B = 0
C = 0

SI (B == 1) AND (A / C > 5) ALORS
    AFFICHER "Vrai"
SINON
    AFFICHER "Faux"
FIN SI

\end{lstlisting}
\end{Exercise}
% ### **Exercice 3 : Application Séquentielle**

% **Réponse (Masquée) :**

% 1.  Le programme commence par évaluer la première partie de l'expression $AND$ : $(B == 1)$.
% 2.  Puisque $B = 0$, l'expression $(B == 1)$ est **Faux ($0$)**.
% 3.  L'opérateur $AND$ utilise l'évaluation séquentielle (court-circuit). Dès que la première partie est Faux, il sait que le résultat final sera Faux, quel que soit le second terme.
% 4.  Le programme **ne tente jamais** d'évaluer la seconde partie $(A / C > 5)$, évitant ainsi la division par zéro ($A/0$).
% 5.  Comme l'expression complète est jugée Faux, le bloc `SINON` est exécuté.

% **Résultat affiché :** "Faux"